\input{preamble}

\title{Section 3.9 --- Equations with Rational Exponents --- Answer Key}

\begin{document}
\maketitle

\section*{Equations with Rational Exponents}
Solve each equation. If there is no solution, write ``no solution.'' Give irrational solutions as simplified radicals. You can use your calculator to evaluate rational powers.
\begin{smenum}
\mitemxxx
{$x=32$}
{$x=16$}
{$x=-27,27$}
\mitemxxx
{$x=-3125$}
{$x=-\sqrt[7]{8}$}
{No solution}
\mitemxxx
{No solution}
{$x=\sqrt[3]{25}$}
{$x=\dfrac{\sqrt[7]{432}}{2}$}
\mitemxxx
{$x=\frac{33}{2}$}
{$x=3-\sqrt[4]{1000},3+\sqrt[4]{1000}$}
{$x=\frac{11}{4}$}
\end{smenum}

\section*{Pretend Polynomials}
Solve each equation by treating it as a polynomial. You might be able to solve by factoring, or you might need to use the quadratic formula. If solutions are irrational, give a decimal approximation to two decimal places, but don't round intermediate steps.
\begin{smenum}
\mitemxx
{$x=-64,2$}
{$x=9,25$}
\mitemxx
{$x=-1,\frac{1}{6}$}
{$x=-\frac{2}{3},0,\frac{5}{2}$}
\mitemxx
{$x=-\frac{8}{27},\frac{1}{512}$}
{No solution}
\mitemxx
{$x\approx 2.88,28.12$}
{$x\approx -30.75,30.75$}
\mitemxx
{$x=-27,27,\frac{1}{8}$}
{This problem has a typo which prevents the polynomial from being factored; skip it.}
\end{smenum}

\end{document}
